\documentclass[10pt]{article}
\usepackage[margin=20mm]{geometry}
\usepackage{amsmath,amssymb,mathtools,fontspec,unicode-math,xurl,hyperref}
\IfFontExistsTF{DejaVu Serif}{\setmainfont{DejaVu Serif}}{\setmainfont{Latin Modern Roman}}
\setmathfont{latinmodern-math.otf}
\setlength{\parindent}{0pt}\setlength{\parskip}{6pt}\sloppy
\title{A verification of selected arguments in
Local Integral Transforms and Global Spectral Decomposition}
\author{}
\date{}
\begin{document}
\maketitle
Abstract: The acquired author manuscript is examined with attention to local normalizations, Mellin contours, the passage to the global spectrum, and the short-window application. An explicit unramified example contradicts the holomorphy and absolute-convergence strip stated in Theorem 2(1). The inversion statement also needs to exclude its diagonal argument or specify a limiting interpretation. Selected global normalization and residue calculations are consistent. These are initial verification findings; scientific adjudication and independent acceptance remain pending.\par
\section*{Identity and scope}
The source is Valentin Blomer, Subhajit Jana and Paul D. Nelson, Local integral transforms and global spectral decomposition, arXiv:2404.10692v2, 22 May 2025, 41 physical pages [BJN]. The assigned journal identifier is DOI 10.1007/s00039-025-00714-0. The report concerns these exact acquired preprint bytes; equivalence with the 57-page journal version has not been established. Page references below use the acquired PDF’s physical numbering, which agrees with its printed pages.\par
The checks concern Theorems 1–5 and selected supporting arguments. This is not a new proof of the full Whittaker–Plancherel theorem, Ichino’s formula, or the cited spectral moment estimates. Those inputs, and the unverified analytic estimates identified in the ledger, remain dependencies.\par
\section*{Setting and checked distinctions}
The local field is of characteristic zero (§2.1), with self-dual additive measure \(dy\) and \(d^\times y=dy/|y|\). The normalized principal series is denoted \(I(\chi)\), and its inducing character satisfies \(f(n(x)a(y)g)=\chi(y)|y|^{1/2}f(g)\). The introduction permits \(\pi_2=I(\chi_2)\) for some inducing character; it imposes no sign condition on \(\Re\chi_2\).\par
For the forward transform in Theorem 1, the strip is \[\vartheta_1+\vartheta_2<\sigma<\tfrac12-\vartheta.\] Its inverse is stated only for \(y_1\ne y_2\). The paper explicitly distinguishes absolute convergence of the indicated iterated inverse integrals from joint absolute convergence (pp. 4–5). The latter must not be inserted as an unstated Fubini hypothesis. The change from a shifted pair \((z,yz)\) to the Mellin-over-shift transform in §2.7 is useful, but it does not by itself remove poles of an individual gamma factor.\par
\section*{Problem 1: the strip in Theorem 2(1)}
Assessment: Major, stated local transform theorem. Theorem 2(1), p. 5, claims holomorphy and absolute convergence for \(\vartheta_1<\Re\chi<\tfrac12-\vartheta\), omitting \(\vartheta_2\). Its proof on p. 17 integrates Theorem 1(1). The larger stated strip admits the following counterexample.\par
\subsection*{Admissible data and noncancellation}
Take \(F=\mathbb Q_3\), an unramified additive character, and self-dual measure normalized by \(\operatorname{vol}(\mathbb Z_3)=1\). Set \[\pi_1=\pi=I(1),\qquad
 \pi_2=I(|\cdot|^{-1/4}),\qquad
 \chi_2=|\cdot|^{-1/4},\qquad \chi_0=1.\] These are irreducible generic unitary principal series, with \(\vartheta_1=\vartheta=0\) and \(\vartheta_2=1/4\); the second is complementary series. Thus even the stronger sum condition of Theorem 1 holds. Choose compact smooth Kirillov restrictions \[W_1(a(z))=\mathbf 1_{\,1+3\mathbb Z_3}(z),\qquad
 W_2(a(z))=\mathbf 1_{\,2+3\mathbb Z_3}(z).\] The compact Kirillov model allows these test vectors. Their values are real, so complex conjugation of \(W_2\) changes nothing.\par
For \(h(z,yz)=W_1(a(z))\overline{W_2(a(yz))}\), definition (1.10) gives \[H(y,1)=\int_{\mathbb Q_3^\times}h(z,yz)\,d^\times z
       =\frac13\,\mathbf 1_{\,2+3\mathbb Z_3}(y).\] Indeed \(z\equiv1\pmod3\), and \(yz\equiv2\pmod3\) is equivalent to \(y\equiv2\pmod3\). On this support both \(|y|\) and \(|1-y|\) equal one. For the unramified character \(\chi=|\cdot|^s\), the inner integral in (1.12) is therefore exactly \(1/9\), independently of \(s\). In particular, neither the test weight nor the factor involving \(1-y\) can cancel a gamma-factor pole.\par
Equation reference: eq:pole. For the chosen unramified data, definition (1.6) gives \[\gamma(u,I(1))
 =\left(\frac{1-3^{-u}}{1-3^{u-1}}\right)^2.\] The full \(\chi\)-integrand in Theorem 2(1), on its unramified character component, is consequently \[
 \frac19
 \left(\frac{1-3^{-1/2-s}}{1-3^{-1/2+s}}\right)^2
 \left(\frac{1-3^{-5/4+s}}{1-3^{1/4-s}}\right)^2 .\] At \(s=1/4\), the last denominator has a double zero; both numerators and the first denominator are nonzero. This is a genuine double pole. It lies strictly inside the advertised strip \(0<\Re s<1/2\). On the permitted contour \(\Re s=1/4\), the absolute value is comparable to \(|\Im s|^{-2}\) near zero, and the character integral is not absolutely convergent. The left-hand transform with \(\chi_0=1\) is well-defined by the stated Whittaker decay estimate.\par
\subsection*{Scope and repair}
This disproves the two assertions about the stated strip, not the existence of the local transform or every identity in the paper. It is a defect in a main local theorem, rather than a harmless choice of normalization. The Major assessment concerns that scope; a short hypothesis correction does not turn the false assertion into a peripheral typographical issue.\par
A conservative correction retains \[\vartheta_1+\vartheta_2<\sigma<\tfrac12-\vartheta,
 \qquad \vartheta_1+\vartheta_2+\vartheta<\tfrac12,\] when deriving (1.12) directly from Theorem 1(1). The example’s pole is then excluded. Any wider strip requires an explicit continuation prescription, residues where necessary, or additional hypotheses on the inducing character. This report does not certify an optimal replacement theorem. In particular, an unstated decision to use the other inducing character cannot repair the theorem as quantified.\par
\section*{Problem 2: the diagonal in Theorem 2(2)}
Assessment: Minor, omitted domain restriction. Theorem 2(2), p. 5, allows \(y\in F^\times\), while its right-hand side contains \(\chi(1-y)\) and division by \(\chi_0(1-y)|1-y|\). At the allowed value \(y=1\), these expressions are undefined. The proof on p. 17 substitutes \(y_1=z,\ y_2=yz\) into Theorem 1(2), whose hypothesis requires \(y_1\ne y_2\), and then changes variables by multiplication by \((1-y)^{-1}\). Thus the proof itself requires \(y\ne1\).\par
For compact smooth \(h\), the defining integral \(H(y,\chi_0)\) is well-defined and smooth at \(y=1\): the \(z\)-support stays in a fixed compact subset of \(F^\times\), so differentiation, or local constancy in the nonarchimedean case, passes under the integral. A precise formulation states the displayed formula for \(y\ne1\), with the diagonal recovered by this limiting value if desired. This corrects the domain without purporting to resolve Problem 1’s contour issue.\par
\section*{Global and application checks}
\subsection*{Normalization and the second diagonal}
The cuspidal pairing on p. 20 is half the Petersson integral, and the cuspidal spectral measure on p. 21 is half counting measure. These two factors cancel in the spectral projector. In the triple-product calculation on pp. 26–27, writing \(Z=\zeta_F^{(S)}(2)\), the explicit scalar product \((Z/8)(8/Z^3)Z^2\) equals one. The remaining adjoint factor is squared, as in Lemma 3.6. This checks the displayed algebra conditional on the cited triple-product and local integral identities.\par
For Theorem 3, the nonzero shift \(b\) kills the one-dimensional spectral contribution: its restriction to the unipotent subgroup is constant, and the nontrivial Fourier coefficient is zero (p. 28). The principal-series assumption at the places in \(S\) is explicit and essential to the chosen local trilinear functional; the source does not claim a general three-supercuspidal formula.\par
For Theorem 4, the two poles in Lemmas 3.11–3.12 contribute with opposite contour signs. Since the corresponding residues are negatives, the \(1/2\) in the Eisenstein measure leaves one full residue, yielding the second diagonal term on p. 34. Omitting that term would be incorrect. The comparison of \(I\) and \(I_0\) is performed by continuation with residues; their identical formal integral expressions alone do not make them equal. The external convergence and continuation estimates are not independently reproved here.\par
\subsection*{The short-window calculation}
With \(n_1=n+b,\ n_2=n\), the weight printed at the start of §4 (p. 34) equals \(V((n+b-X)/Y)\). Thus that displayed sum represents \(S(X-b,Y,b)\), with the definition on p. 8. The required weight for \(S(X,Y,b)\) is \[\sqrt{|n_1n_2|}\,
 V_1\!\left(\frac{n_1-(X+b)}Y\right)
 V_2\!\left(\frac{n_2-X}Y\right).\] This is a local centering correction. Running the subsequent auxiliary calculation with center \(X+b\) preserves the scale: \(X\le X+b\le3X/2\). In the mean-square case the shifted interval \([X+b,2X+b]\) lies in \([X+b,2(X+b)]\). This correction supplies no counterexample to the asserted estimates.\par
The integration-by-parts identity on p. 36 was checked directly. Differentiation in \(t\) produces \(b/Y\) and \(b(1+\tau)/Y\); integration by parts in \(x\) combines them into the stated \(b\tau/Y\) term and \(b/X\) derivative of the outer weight. The support bound \(|\tau|\ll Y/X\) then gives \(b|\tau|/Y\ll b/X\). The final displayed mean-square bound implies (1.18) using \(\sqrt{b/Y}\le(1+b/Y)/2\).\par
A further contour notation needs checking in any revision: (4.11) is printed with initial lines \(\Re s_j=3/4\), whereas the following paragraph discusses residues at \(1/2\) on shifting to \(1+\varepsilon\). The source’s original transform starts to the left of those poles. The complete application estimate, including all residue and spectral moment bounds, is not certified by this report; this observation is retained in coverage and is not counted as an additional proved theorem failure.\par
\section*{Conclusion}
Problem 1 is a supported Major finding about the scope of a main local transform theorem. Problem 2 is a local domain omission. The checked global normalization and residue algebra provide counterevidence to broader claims that the spectral method fails. No statement is made about author adoption or subsequent journal corrections. Scientific decision is pending, there are no human review entries, and independent acceptance is incomplete.\par
\section*{Evidence and limitations}
The pole and test-vector calculations above are analytic; no finite numerical experiment or author code substitutes for the proof. The original source extraction is reused with its historical execution origin. The hypergeometric appendix, all cited external proofs, and every interchange or estimate outside the stated checks are not claimed as independently verified. The absence of optional author code in the bounded acquisition search is not a mathematical defect.\par
1 V. Blomer, S. Jana and P. D. Nelson, Local integral transforms and global spectral decomposition, arXiv:2404.10692v2, 22 May 2025.\par
\end{document}
